\section{Background}
Molecular orbitals are expressed as linear combinations of atomic basis functions. The coefficients, orbital energies, occupation numbers, and basis-set information are normally stored in a molecular orbital file obtained from an external quantum chemistry calculation. The \texttt{postorb} module can read these data independently of an electronic excited-state calculation. A single file is sufficient for most orbital analyses. When an analysis requires orbitals from separate calculations, the \texttt{\$\$mopack} input described in Section~\ref{sec-mofiles} can provide multiple orbital files.

Orbital inspection is often a useful first step when selecting states or assessing the orbitals used in an electronic-structure calculation. The available outputs include the molecular geometry, AO and MO information, selected orbital coefficients, and exported orbital files in supported formats. AO and MO overlap integrals measure the overlap of basis functions and orbitals. Other available one-electron integrals include electric dipole, magnetic dipole, velocity, and selected Cartesian multipole operators. The required orbital indices and property operators are selected in Job Control below.

The module can also generate cube files for selected AO and MO wave functions, electron density, and spin density. A cube file contains values evaluated on a three-dimensional spatial grid and can be displayed using suitable visualization software. The electron-density and spin-density calculations use the orbital occupation information in the supplied orbital file. These orbital-based densities should be distinguished from state-specific densities and density differences calculated from electronic-state wave functions with \texttt{postwfn}. A larger requested spatial grid may improve a visualization but also increases computational effort and the size of the output.
